% ID: SEL-D0025-OB-E03
% ESTADO: cotejada. Segunda pasada visual completa y borrador compilado revisados.
% Fuente: Selectividad/matematicas_aplicadas/suplemente_junio_EXAMEN 4-ANDALUCÍA.pdf
% SHA-256: c1a3a7a37d38b04359115c64762b1f8a543e4d66155338931de84ef6c304da2c
% Página: 2; opción B; ejercicio original 3.

Los alumnos que cursan una asignatura deben realizar dos exámenes: uno teórico y otro práctico. El 50\% de los alumnos aprueba los dos exámenes, el 6\% no aprueba ninguno y el 20\% solo aprueba el teórico. Se elige un alumno al azar.
\begin{enumerate}[label=\alph*)]
\item \textbf{(1 punto)} \textquestiondown{}Cuál es la probabilidad de que apruebe al menos uno de los dos exámenes?
\item \textbf{(1.5 puntos)} Si ha aprobado el teórico, \textquestiondown{}cuál es la probabilidad de que no apruebe el examen práctico?
\end{enumerate}
